% ===================== 第 5 章 分题型解题策略 =====================
\section{分题型解题策略}

\subsection{写作：四类框架与模板}

\begin{tipbox}[title=时间分配（30 分钟）]
审题列提纲 5 分钟 $\to$ 写作 20 分钟 $\to$ 检查 5 分钟。写作在听力之前进行，写完后不要再改，把精力留给听力。
\end{tipbox}

\subsubsection*{1. 现象解释型}

\textbf{适用}：题目给出一种社会现象（如沉迷虚拟世界、追星、网络虚假信息）。

\begin{itemize}
  \item 首段：引出话题 + 点明现象 + 表明态度（2--3 句）
  \item 主体：现象的表现 $\to$ 产生的原因（2 条）$\to$ 影响（1 条）
  \item 结尾：总结观点 + 提出建议或呼吁
\end{itemize}

模板句：
\begin{itemize}
  \item Recent years have witnessed a growing trend of \ldots
  \item Several factors can account for this phenomenon. To begin with, \ldots\ In addition, \ldots
  \item Only by \ldots\ can we \ldots
\end{itemize}

\subsubsection*{2. 问题解决型}

\textbf{适用}：题目要求分析问题并给出措施（如如何平衡工作与生活、如何应对创新不足）。

\begin{itemize}
  \item 首段：描述问题及其严重性
  \item 主体：问题产生的原因（简写）$\to$ 两到三条具体措施（重点）
  \item 结尾：预期效果 + 号召
\end{itemize}

模板句：
\begin{itemize}
  \item It is high time that we took effective measures to \ldots
  \item First and foremost, the government / universities / individuals should \ldots
  \item What's more, it is advisable to \ldots
  \item Only when all sides make joint efforts can the problem be solved.
\end{itemize}

\subsubsection*{3. 观点选择型}

\textbf{适用}：题目要求比较两种做法并给出选择（如主修文科还是理科）。

\begin{itemize}
  \item 首段：点明争论话题 + 表明自己的选择
  \item 主体：让步（对方观点的合理性，1 句）$\to$ 己方理由（2--3 条，配例证）
  \item 结尾：重申观点
\end{itemize}

模板句：
\begin{itemize}
  \item When it comes to \ldots, opinions vary from person to person.
  \item Admittedly, \ldots\ However, I am firmly in favor of \ldots
  \item From my perspective, the advantages far outweigh the disadvantages.
\end{itemize}

\subsubsection*{4. 谚语评述型}

\textbf{适用}：题目给出一句名言或谚语，要求评论（如“明日之最佳准备，就是把今天做到最好”）。

\begin{itemize}
  \item 首段：复述并解释名言的含义 + 点明观点
  \item 主体：道理阐述（1 条）$\to$ 举例论证（1--2 个例子，可用个人经历、名人故事、社会现象）
  \item 结尾：联系自身，说明如何践行
\end{itemize}

模板句：
\begin{itemize}
  \item As the saying goes, \ldots\ The message conveyed is that \ldots
  \item A case in point is that \ldots
  \item To sum up, the saying reminds us that \ldots
\end{itemize}

\begin{notebox}[title=避坑清单]
\begin{itemize}
  \item 首句若已给出，必须照抄，不得改写、翻译或另起一句。
  \item 字数控制在 150--200 词，超过或不足都会扣分；段落不要少于三段。
  \item 避免三种低级错误：主谓一致、时态混用、名词单复数。
  \item 不要通篇堆砌 “I think”；用 From my perspective / It is my belief that 等替换。
\end{itemize}
\end{notebox}

\subsection{翻译：五步流程与五类技巧}

\subsubsection*{五步流程（建议 25--30 分钟）}

\begin{enumerate}
  \item \textbf{通读}全文，判断话题类别（文化 / 社会 / 地理），圈出专有名词与难点词。
  \item \textbf{拆句}：把中文长句按“意群”切成短句，标出主从句关系。
  \item \textbf{直译}主干：先保证“主谓宾”正确，再补修饰成分。
  \item \textbf{升级}表达：替换低级词、调整语序、适当合句、必要时改被动。
  \item \textbf{检查}：时态、单复数、冠词、拼写、漏译。
\end{enumerate}

\subsubsection*{五类技巧}

\begin{center}
\small
\begin{longtable}{p{3.1cm}p{11.5cm}}
\toprule
\textbf{技巧} & \textbf{示例} \\
\midrule
\endfirsthead
\toprule
\textbf{技巧} & \textbf{示例} \\
\midrule
\endhead
词类转换 & 他很喜欢读书 $\to$ He has a great love for reading.（动词转动词性名词） \\
\midrule
语序调整 & 他昨天在公园遇见了朋友 $\to$ He met his friend in the park yesterday. \\
\midrule
语态转换 & 这部小说被翻译成多种语言 $\to$ The novel has been translated into many languages. \\
\midrule
拆分合并 & 中文多短句 $\to$ 英语用从句、分词、同位语合并 \\
\midrule
固定表达 & “随着……的发展” $\to$ with the development of；“人们普遍认为” $\to$ It is widely believed that \ldots \\
\bottomrule
\end{longtable}
\end{center}

\begin{tipbox}[title=翻译高频“官方译法”提醒]
\begin{itemize}
  \item 二十四节气 $\to$ the 24 Solar Terms；国粹 $\to$ the quintessence of Chinese culture（不要译作 national treasure）。
  \item 朝代一律音译 + Dynasty：the Tang Dynasty。
  \item “位于” $\to$ be located in / lie in；“流经” $\to$ flow through；“追溯到” $\to$ date back to。
  \item 数字与年代：公元前 $\to$ BC；公元后 $\to$ AD；“亿” $\to$ hundred million；注意 \emph{billion} = 十亿。
\end{itemize}
\end{tipbox}

\subsection{阅读理解：分题型打法}

\begin{enumerate}
  \item \textbf{选词填空（建议 8--10 分钟）}：先给 15 个选项标词性、分四组；再逐空判词性、看搭配、比词义；最后通读验证。记不住的不纠结，一空 3.55 分，性价比最低，可放到最后做。
  \item \textbf{长篇阅读（建议 15 分钟）}：先读题干圈关键词（人名、数字、专有名词），再按段落扫读定位，重点关注\textbf{同义替换}；一题一勾，防止重复使用。
  \item \textbf{仔细阅读（建议 30--35 分钟）}：每篇 8--10 分钟，先题后文、定位作答；主旨题看首尾段，态度题看形容词，推理题注意“同义改写即答案”。
\end{enumerate}

\subsection{听力：三个动作}

\begin{enumerate}
  \item \textbf{预读}：利用题目间隙读选项，圈出选项差异词（数字、地点、态度）。
  \item \textbf{抓信号词}：but, however, in fact, actually, the most important thing is --- 转折与强调之后往往是答案。
  \item \textbf{边听边涂}：听力一结束就收答题卡 1，答案必须随时填涂。
\end{enumerate}

\subsection{考场时间总方案}

\begin{center}
\small
\begin{longtable}{p{3.0cm}p{5.2cm}p{6.4cm}}
\toprule
\textbf{时段} & \textbf{任务} & \textbf{要点} \\
\midrule
\endfirsthead
\toprule
\textbf{时段} & \textbf{任务} & \textbf{要点} \\
\midrule
\endhead
15:10--15:40 & 写作 & 30 分钟一气呵成，留 5 分钟检查 \\
\midrule
15:40--16:10 & 听力 & 边听边涂卡，S1 结束后不再翻看 \\
\midrule
16:15--16:40 & 仔细阅读（前 4--5 题） & 先做分值最高的部分，进入状态 \\
\midrule
16:40--17:05 & 翻译 & 预留 25 分钟；写完直接誊到答题卡 2 \\
\midrule
17:05--17:20 & 长篇阅读 & 定位 + 同义替换 \\
\midrule
17:20--17:25 & 选词填空 & 填满不留空；全部撞一遍核对是否有重复用词 \\
\bottomrule
\end{longtable}
\end{center}
